\medskip

\begin{enumerate}
	\item \textbf{Affirmation 1 : FAUSSE.}

	La propriété du cours indique que toute suite décroissante et minorée par 0 converge vers une limite $\ell$, avec $\ell\geqslant 0$.

	Il suffit donc d'exhiber un contre-exemple.

	La suite constante égale à 1 est décroissante (mais pas strictement décroissante) et minorée par 0 (entre autres), et pourtant, elle converge vers 1, et pas vers 0.

	Si on préfère donner un contre exemple avec une suite strictement décroissante, on peut par exemple choisir la suite définie sur $\N^*$ et de terme général :\quad $u_n = 1 + \dfrac{1}{n}$, par exemple. Cette suite est assez clairement décroissante, minorée par 0 et converge vers 1.

	\item \textbf{Affirmation 2 : VRAIE.}

	En effet, pour $n$ entier naturel :\quad $v_n = \dfrac{-9^n + 3^n}{7^n} = -\left(\dfrac{9}{7}\right)^n \times \dfrac{1 - \left(\dfrac{3}{9}\right)^n }{1} = -\left(\dfrac{9}{7}\right)^n \times \left(1 - \left(\dfrac{1}{3}\right)^n\right)$

	Comme on a : $\dfrac{9}{7}>1$, on en déduit que $\lim\limits_{n \to +\infty} \left(\dfrac{9}{7}\right)^n = +\infty.$

	Par ailleurs : $-1 < \dfrac{1}{3}<1$, donc on en déduit $\lim\limits_{n \to +\infty} \left(\dfrac{1}{3}\right)^n = 0$, donc, par limite de la somme : $\lim\limits_{n \to +\infty} 1 - \left(\dfrac{1}{3}\right)^n = 1$.

	Finalement, par limite du produit : $\lim\limits_{n \to +\infty} -\left(\dfrac{9}{7}\right)^n \times \left(1 - \left(\dfrac{1}{3}\right)^n\right) = -\infty$.

	Si on a, pour tout $n$ naturel, $u_n \leqslant v_n$ et $\lim\limits_{n \to +\infty} v_n = -\infty$, alors, par comparaison : \quad $\lim\limits_{n \to + \infty} u_n = -\infty$.

	\item \textbf{Affirmation 3 : VRAIE.}

	L'appel \texttt{terme(4)} commence par initialiser la variable \texttt{U} avec la valeur 1 (ligne 2 de la fonction).

	Puis, la boucle \texttt{for} va s'exécuter \texttt{N} fois, ici donc 4 fois, avec le compteur \texttt{i} qui va prendre les valeurs entières entre 0 et \texttt{N-1}, donc ici 0, 1, 2 et 3.

	\begin{itemize}
		\item La première exécution \og modifie\fg{} \texttt{U} en \texttt{U + 0}, la valeur reste égale à 1;
		\item La deuxième exécution modifie \texttt{U} en \texttt{U + 1}, la valeur devient égale à $1 + 1 = 2$;
		\item La troisième exécution modifie \texttt{U} en \texttt{U + 2}, la valeur devient égale à $2 + 2 = 4$;
		\item La dernière exécution modifie \texttt{U} en \texttt{U + 3}, la valeur devient égale à $4 + 3 = 7$.
	\end{itemize}

	On a donc bien la valeur 7 renvoyée par cet appel.
	\item \textbf{Affirmation 4 : FAUSSE.}

	Pour connaître le montant total du prix A, il suffit de multipler 1000 par 15. Le montant total est donc de $\np{15000}$\,\euro{}.

	Pour le prix B, il s'agit d'additionner les 15 premier termes d'une suite géométrique de premier terme 1 (le montant du prix le premier jours) et de raison 2 (car le montant chaque jour est le double du montant de la veille).

	On applique la formule connue : \quad $1 \times \dfrac{1 - 2^{15}}{1 - 2} =\dfrac{1-2^{15}}{-1} = 2^{15} - 1 = \np{32767}$.

	Comme $\np{32767} > \np{15000}$, le prix B est (nettement) plus avantageux.

	\emph{Remarque :} En cas d'oubli de cette formule, on peut aussi (patiemment) calculer la somme des 15 termes, voire même remarquer que la somme reçue au quinzième jour sera de $2^{15 - 1} = \np{16384}$, donc rien que la somme reçue le quinzième jour du prix B est strictement supérieure aux quinze jours cumulés pour le prix A.

	\item \textbf{Affirmation 5 : VRAIE.}

	Soit $n$ un entier naturel non nul quelconque :

	\begin{itemize}
		\item La fonction ln est strictement croissante sur $\R^{*+}$, et $\ln(1) = 0$, donc on en déduit que la fonction ln est à valeurs positives sur $[1 ~;~ +\infty[$ (voire à valeurs strictement positives sur $]1~;~ +\infty[$).
		\item $\begin{aligned}[t] \displaystyle v_{n+1}	- v_n &= \int_{1}^{n+1} \ln x \:\mathrm{d} x - \int_{1}^{n} \ln x \:\mathrm{d} x \\
		&= \int_{n}^{n+1} \ln x \:\mathrm{d} x \quad\text{par la relation de Chasles} \end{aligned}$

		\item l'expression $v_{n+1} - v_n$ est donc l'intégrale, entre deux bornes ordonnées dans l'ordre croissant (car $n < n+1$) d'une fonction à valeurs positives sur l'intervalle d'intégration (car $[n ~;~ n+1] \subset [1 ~;~ +\infty[$, puisque $n\geqslant 1$) : cette expression est donc positive.

\end{itemize}
	Pour un $n$ quelconque supérieur à 1, la différence entre les termes $v_{n+1}$ et $v_n$ est positive, donc la suite $\left(v_n\right)$ est bien croissante.
\end{enumerate}
